%\documentclass[12pt, a4paper]{abnt}
\documentclass[12pt, brazil, ruledheader, pnumromarab,normaltoc]{abnt}
%\documentclass[brazil, ruledheader]{abnt}
\usepackage[utf8]{inputenc}
%\usepackage[T1]{fontenc}
%\usepackage{dsfont}
%\usepackage{verbatim}

%\chaptersize{\large}
%\fontfamily{cmr}\fontseries{b}\selectfont
%\renewcommand{\ABNTchapterfont}{\bfseries\sffamily\fontseries{b}\selectfont}
%\renewcommand{\ABNTchapterfont}{\bfseries\fontfamily{cmr}\fontseries{b}\selectfont} % Computer Modern Bold
%\renewcommand{\ABNTchapterfont}{\bfseries\fontfamily{cmss}\fontseries{sbc}\selectfont} % CM Sans Serif Semibold Condensed
\renewcommand{\ABNTchapterfont}{\bfseries\fontfamily{ptm}\fontseries{b}\selectfont} % Times Roman Bold
%\renewcommand{\ABNTsectionfont}{\bfseries\fontfamily{cmr}\fontseries{b}\selectfont}
\renewcommand{\ABNTsectionfont}{\bfseries\sffamily}
%\renewcommand{\ABNTchapterfont}{\upshape}

\makeatletter
\usepackage[brazil]{babel}
\makeatother

%\usepackage[num, abnt-alf, abnt-verbatim-entry=yes, cite]{abntcite}
%\usepackage[abnt-alf]{abntcite}
\usepackage{graphicx, color}
%\usepackage{url}
%\usepackage{makeidx}
%\usepackage{amsfonts}
%\usepackage{amssymb}
%\usepackage{indentfirst}
\usepackage[citecolor=red,colorlinks=true,linkcolor=blue,hyperindex, pdfborder={0 0 0}] {hyperref}

%\usepackage{a4wide}            % correta formatação da página em A4
%\usepackage{setspace} 
%\usepackage{fancyhdr}
\usepackage[normalem]{ulem}
%\pagenumbering{roman}

% Inserindo as informações do documento{author, title}
\autor{{André Luiz Ramos de Souza} \\ {Diego de Souza Lopes}}

%\titulo{Controle de Acesso à Rede com PacketFence}
\titulo{Controle de Acesso à Rede\vspace{0.2cm}\\com\vspace{0.2cm}\\PacketFence}

\orientador{Prof. MsC. Maurício Severich}

\instituicao{FACULDADE DE TECNOLOGIA SENAI DE  DESENVOLVIMENTO GERENCIAL - FATESG}

\comentario{Monografia apresentada a Faculdade de Tecnologia SENAI de Desenvolvimento Gerencial como requisito a obtenção do título de Tecnólogo em Redes de Computadores.}

\local{GOIÂNIA}

\data{2011}

\begin{document}

% --------------------------- DESENHANDO A CAPA --------------------- %
\begin{titlepage}
\begin{center}

%\vspace{0.1\textheight}

\includegraphics{senai}

%\vspace{0.1cm}

\textbf{FACULDADE DE TECNOLOGIA SENAI DE\\DESENVOLVIMENTO GERENCIAL - FATESG}

\textsc{SUPERIOR EM TECNOLOGIA EM REDES DE COMPUTADORES}

%\hspace{.45\textwidth}
\vspace{3cm}

\large André Luiz Ramos de Souza\\Diego de Souza Lopes

%\hspace{.35\textwidth}
\vspace{3cm}

\textbf{\LARGE Controle de Acesso à Rede\vspace{0.1cm}\\com\vspace{0.3cm}\\PacketFence}

%\hspace{.35\textwidth}
\vspace{2.5cm}

%\textbf{Orientador:} MsC. Maurício Severich

%\hspace{.35\textwidth}
\vspace{3.5cm}
%\vfill

\normalsize GOIÂNIA\\

%\vspace{.1\textwidth}
%\large 2011
\normalsize 2011

\end{center}
\end{titlepage}
% --------------------------- FIM CAPA ------------------------------- %
\folhaderosto
%\folhadeaprovacao

% --------------- FICHA CATALOGRAFICA ----------------------- %
\begin{minipage}{7cm}
\begin{center}
%{\normalsize \bf \MakeUppercase{Ficha catalográfica}}
\vspace{.15\textwidth}
\begin{tabular}{|cl|} \hline
  \hspace{1.3cm} & \\
  & Ramos, André Luiz de Souza\\
  & \\
  \hspace{0.2cm}  & \hspace{0.3cm} Controle de Acesso à Rede com PacketFence / André Luiz\\
  &  Ramos de Souza, Diego de Souza Lopes. - Goiânia: 2011.\\
%  & FATESG, 2011. \\
  & \hspace{0.65cm} 58f. : il. \\
  & \\
  & \hspace{0.6cm} Orientador: Prof. MsC. Maurício Severich.\\
  & \\
  & \hspace{0.6cm} Trabalho de conclusão de curso (Redes de Computadores) -\\
  & \hspace{0.6cm} Faculdade de Tecnologia SENAI DE DESENVOLVIMENTO \\
  & \hspace{0.6cm} GERENCIAL, Goiânia, 2011.\\
  & \\
  & \hspace{0.6cm} Inclui bibliografia. \\
% & \hspace{0.6cm} 1 - Detecção e classificação de faltas \hspace{0.1cm} 2 - 
%Linhas de transmissão \\
% & \hspace{0.6cm} 3 - Redes neurais artificiais \hspace{0.1cm} 4 - 
%Transformada wavelet \hspace{0.1cm} 5 - Título\\
  & \\
% & \hspace{9.75cm} CDU 631.317.35 \\
 \hline
\end{tabular}
\end{center}
\end{minipage}
% --------------------------- FIM ------------------------------- %


% --------------------------- FOLHA APROVACAO ---------------------- %
\begin{folhadeaprovacao}
\setlength{\ABNTsignthickness}{0.4pt}
\setlength{\ABNTsignskip}{1cm}
\vspace*{0.2cm}
\begin{center}
\large André Luiz Ramos de Souza\\Diego de Souza Lopes\\
%\hspace*{2cm}
%\end{center}
\vspace*{1.5cm}
\textbf{\LARGE Controle de Acesso à Rede\vspace{0.1cm}\\com\vspace{0.3cm}\\PacketFence}
\end{center}
\vspace*{1cm}
Monografia apresentada a Faculdade de Tecnologia SENAI de Desenvolvimento Gerencial como requisito a obtenção do título de Tecnólogo em Redes de Computadores.\\
\vspace*{0.2cm}\\
Goiânia, \uline{10}  de \uline{Dezembro} de 2011.\\
%\vspace*{1cm}
\begin{center}
BANCA EXAMINADORA
\end{center}
%\hspace*{1cm}
\assinatura{\textbf{Prof. MsC. Mauricio Severich}\\Orientador}
%\hspace*{1cm}
\assinatura{\textbf{Prof. MsC. Rafael de Tal} \\ Faculdade de Tecnologia SENAI - FATESG}
%\hspace*{1cm}     
\assinatura{\textbf{Prof. MsC. Diogo de Tal} \\ Faculdade de Tecnologia SENAI - FATESG} 
%\hspace*{1cm}
%\assinatura{\textbf{Prof. MsC. Fulano de Tal} \\ Faculdade de Tecnologia SENAI - FATESG} 
      
\end{folhadeaprovacao}
% ----------------------- FIM FOLHA APROVACAO --------------------- %
\sumario
%\folhaderosto
%\folhadeaprovacao
%\tableofcontents
%\listoffigures  % Lista de Figuras
%\listoftables  % Lista de Tabelas

% ------------------------------- FIM ------------------------------- %

% ------------------------------ RESUMO ----------------------------- %
\begin{resumo}
%$\phantom{linha em branco}$\\

Um aspecto pouco observado e que traz grandes problemas para uma organização é a falta de controle de acesso em sua infraestrutrura, desde acesso a rede cabeada à sem fio. Atualmente os pontos de redes que estão em desuso em uma organização, poderá ser utilizado por pessoas com más intenções \emph{"atacante"}, pois desta forma é possível obter algumas informações básicas através deste ponto de rede, tais como:
\begin{enumerate}
\item[-]{Endereço IP};
\item[-]{Máscara de Rede};
\item[-]{Gateway};
\item[-]{DNS's}
\end{enumerate}
\par Através das informações acima é um bom ponto de partida para um \emph{"atacante"}, explorar-se dessa falha de segurança da organização, pois é possível realizar interceptação de dados (captura), bem como realizar ataques aos recursos da organização (Servidores, Desktops), podendo conseguir dados sigilosos da mesma. Este trabalho explana o controle de acesso à rede, fazendo uso do software \textit{open source\footnote{Termo em inglês criado pela OSI (Open Source Initiative) e refere-se a software também conhecido como software livre}} \textbf{PacketFence}.\\ Será implementado em ambiente real (corporativo) o PacketFence, a organização participante deste trabalho é uma empresa do Grupo José Alves, referimos à Faculdade ALFA.\\

\textbf{Palavras-Chave:} PacketFence, 802.1X, 802.1Q, EAP, Portal Captive, FreeRadius, Snort, Nessus, SNMP, DHCP, DNS.

\end{resumo}
% ---------------------------- FIM RESUMO ---------------------------- %

% ----------------------------- ABSTRACT ----------------------------- %
\begin{abstract}
%$\phantom{linha em branco}$\\
A little noticed aspect and that brings big problems for an organization is the lack of access control in your infraestrutrura, since access to the wired network wireless. Currently the points of networks that are no longer used in an organization, can be used by people with bad intentions \emph {"attacker"}, because this way you can get some basic information through this network point, such as:
\begin {enumerate}
\item [-] {IP Address};
\item [-] {Netmask};
\item [-] {Gateway};
\item [-] {DNS}
\end{enumerate}
\par
Using the information above is a good starting point for a \emph {"attacker"}, to exploit this security flaw in the organization, it is possible to intercept data (capture) as well as conducting attacks on the organization's resources (servers, Desktops), you can achieve the same sensitive data. This paper explains the access control network, using software \textit {open source \footnote {English term created by the OSI (Open Source Initiative) and refers to software also known as free software}} \textbf {PacketFence}.  Will be implemented in a real environment (corporate) the PacketFence, the participating organization of this work is a company of the José Alves, refer to the Alpha School.\\

\textbf {Keywords:} PacketFence, 802.1x, 802.1Q, EAP, Captive Portal, FreeRadius, Snort, Nessus, SNMP, DHCP, DNS.
\end{abstract} 
% ---------------------------- FIM ABSTRACT -------------------------- %

\chapter{INTRODUÇÃO}
Um aspecto pouco observado pelas organizações é a falta de controle de acesso à infraestrutura de rede, principalmente os pontos de rede não utilizados e que continuam ativos em sua infraestrutura, dando total acesso a pessoas com intenções adversas, podendo explorar-se dessa falha para obter dados sigilosos ou realizar ataques aos servidores da organização.
\par
Afim de preencher esse mercado em crescimento contínuo que demanda a cada dia maior controle de acesso a infraestrutura de rede e que surgem novas ameaças e vulnerabilidades constantemente, existe uma ferramenta com custo zero, o \textbf{PacketFence}, o mesmo será testado e implementado. Essa é uma solução denominada \textit{Network Access Control}  ou simplesmente \emph{NAC}.
\par
A ferramenta controlará os acessos dos equipamentos na infraestrutura de redes de computadores, exigindo uma autenticação e registro do dispositivo conectado, além de verificar se no dispositivo conectado existe alguma falha de segurança ou atividade maliciosa no mesmo. Se a resposta for positiva, o mesmo será direcionado para uma área isolada para que sejam realizadas as possíveis correções dos problemas. Posteriormente à correção, esse equipamento (computador ou notebook) conseguirá realizar qualquer atividade na empresa, seja acesso à Internet ou acesso aos sistemas da organização.

%\part{Conceitos Básicos:\\Teoria}
\chapter{CONCEITOS BÁSICOS}
\section{O que é Network Access Control ?}\label{sec:nac}
\emph{Network Access Control} ou simplesmente \emph{NAC} é uma abordagem unificada afim de prover a segurança em uma rede de computadores. De acordo com \cite{interop-nac} NAC é um controle genérico de acesso à rede. Tal controle de acesso, poderia ser implementado simplesmente com o uso de políticas de acesso ou definindo regras no firewall.
\par
Firewall é um dispositivo de rede que tem por objetivo aplicar políticas de segurança, verificando informações da camada de enlace (camada 2 do modelo OSI\footnote{Open Systems Interconnection}) e da camada de rede (camada 3 do modelo OSI). Aplicando políticas de filtragem  de pacotes TCP/IP  e portas.
\par
NAC é um sistema de controle diferente, tal como firewall, pois possuí ferramentas disponíveis para fornecer o controle de acesso focado no usuário \cite{interop-nac}. Ao contrário do firewall, onde esse controle é baseado em regras. Essas regras são baseadas em origem e destino dos pacotes, definindo-se o que é bloqueado e o que é liberado. Diferentemente de uma solução NAC, onde verifica-se a camada de aplicação (camada 7 do modelo OSI).

\subsection{NAC 1.0: Porque falhou}
As organizações estão cada vez mais querendo controlar o acesso a rede com o uso de tecnologias que melhor proteja suas redes e dados.
De acordo com \cite{mfrizzi} a geração NAC 1.0 falhou devido ao foco em bloqueio de clientes e a falta de agilidade, pois o ambiente de novas ameaças surgiam constantemente e novas vulnerabilidades aparecendo a cada dia, surgindo desta forma um problema muito maior em termos de segurança de uma organização.
\par
A patente \cite{nac} implementa dois métodos, o IEEE 802.1x \cite{802.1x} e o método DHCP\footnote{Dynamic Host Configuration Protocol} \ref{sec:dhcp}, e controle na camada 2.

\subsection{NAC 2.0: Um novo modelo para um futuro mais seguro}
Em uma solução completa de NAC \cite{nac} implementa controle de acesso, autenticação IEEE 802.1x veja a seção \ref{sec:802.1x}, captive portal seção \ref{sec:captive}, detecção de intrusão, gerenciamento na camada 2, scanner de vulnerabilidades. 

\subsection{NAC 3.0: A Evolução}
http://blogs.mcafee.com/enterprise/nac/introducing-the-third-generation-of-nac

\section{O que é Network Access Protection ?}

\section{O que é Network Admission Control ?}

\section{O que é PacketFence ?}
PacketFence \cite{url-pf} é um software de código aberto e construido sobre padrões abertos, ou seja, para evitar o fechamento do padrão por parte do fabricante. A seguir a relação de padrões utilizados:
\begin{enumerate}
\item [•] 802.1X seção \ref{sec:802.1x}
\item [•] Simple Network Management Protocol - SNMP seção \ref{sec:snmp}
\item [•] Standard SNMP management information base (MIB) like BRIDGE-MIB, Q-BRIDGE-MIB, IF-MIB, IEEE8021-PAE-MIB
\item [•] RADIUS seção \ref{sec:radius}
\item [•] NetFlow / IPFIX seção \ref{sec:netflow}
\item [•] Wireless ISP Roaming (WISPR)
\end{enumerate}

Atualmente desenvolvido por programadores que estão localizados nos Estados Unidos em sua maioria. Algumas características impressionantes, captive portal para registro e correção de dispositivos, gerenciamento centralizado com fio e sem fio, suporte 802.1x, isolamento de dispositivos problemáticos na camada 2, integração com o IDS Snort e do scanner de vulnerabilidade Nessus. PacketFence pode ser usado para prover efetivamente redes seguras, de pequena a grandes redes heterogêneas.

\subsection{Desenvolvimento Out-of-Band}

\subsection{Desenvolvimento Inline}

\subsection{Suporte 802.1X}

\subsection{Integração Wireless}

\subsection{Suporte VoIP}


\section{Soluções Comerciais NAC}
https://mosaicsecurity.com/categories/81-network-access-control

\chapter{FERRAMENTAS, PADRÕES, PROTOCOLOS E TECNOLOGIAS}
\section{Dynamic Host Configuration Protocol - DHCP}
DHCP é definido na RFC 2131 \cite{rfc-2131}, o mesmo é um protocolo de configuração dinâmica de host.

\subsection{DHCP Fingerprint}\label{sec:dhcp}
DHCP Fingerprint é um identificador quase único para um determinado sistema operacional específico ou dispositivo. Através desse identificador é possível saber o sistema operacional do cliente. Devido a difusão do protocolo DHCP essa é a maneira mais simples de obter-se informação referente ao cliente. Como definido na RFC 2132, obter a identificação do cliente é possível com o uso da opção "Option Code 61 - Client Identification".
\par
As assinaturas são bastante importantes para ferramentas de analise de rede, tais como sniffers de rede, analisadores de protocolo e soluções NAC. %\ref{sec:nac}
\subsection{Como Funciona}
O cliente DHCP requisita um pedido DHCP nessa requisição ele solicita as opções de DHCP tais como DNS Server, Servidor WINS, gateway padrão, etc. Através da ordem como esse pedido é solicitado essas opções é relativamente único, identificando a versão específica do sistema operacional.
\cite{fingerbank}

%\section{Domain Name System - DNS}

\section{Captive Portal}\label{sec:captive}
Captive Portal é um aplicativo de firewall que captura todos os pacotes do cliente e redireciona o tráfego para uma página WEB, onde o cliente deverá fornecer suas credênciais, tais como usuário, senha ou certificado digital. Ao fornecer as credênciais é verificado se são válidas, caso as credênciais sejam válidas, autoriza-se o acesso a rede.
\cite{captive}

\section{FreeRadius}\label{sec:radius}

\section{IEEE 802.1x}\label{sec:802.1x}
http://standards.ieee.org/about/get/802/802.1.html
http://www.intel.com/support/pt/wireless/wlan/sb/cs-008413.htm

\subsection{Extensible Authentication Protocol - EAP}
http://www.ietf.org/rfc/rfc2284.txt

%\section{Servidor Web Apache}

%\section{Banco de Dados - MySQL}

\section{Simple Network Management Protocol - SNMP} \label{sec:snmp}
Com o crescimento das redes de computadores, de pequenas redes domésticas para grandes redes espalhadas geograficamente e somado a isso vários equipamentos de fabricantes diferentes, a gerência dessas infra-estruturas fica cada vez mais complexas. O protocolo SNMP traz para os administradores de redes a minimização dessas complexidades.
\par
O Simple Network Management Protocol (SNMP) começou a ser desenvolvido no início da década de 80 pelo Internet Engineering Task Force (IETF) com o objetivo de simplificar a gerencia de redes. Com esse protrocolo é possível gerenciar switches, roteadores, softwares, qualquer dispositivo que permita recuperar informações de SNMP.

\subsection{Arquitetura SNMP}
Software de gerência de redes não segue o modelo de \emph{cliente/servidor}, modelo em que somente o cliente faz requisição ao servidor, e sim o conceito de gerente e agente. Nesse modelo as requisições podem ser feitas pelo gerente ao agente (operações GET e SET) ou pelo agente ao gerente (operação TRAP).

\subsection{Gerente}
A aplicação gerente é responsável pelo monitoramento e gerenciamento dos dispositivos gerenciados e também fornecer a interface para o gerente humano. Também conhecido como \emph{Network Management Stations - NMS}, são responsável pelas operações de \emph{pulling} e receber \emph{traps}.
\par
\emph{Pull} é a ação realizada pela NMS para ler dados na base de dados dos agentes e \emph{Trap} é a ação realizada pelos agentes para enviar dados (eventos) dos dispositivos gerenciados para o gerente ou NMS, as \emph{traps} também podem executar algum tipo de ação se a NMS solicitar.
\par Desta forma, não sendo necessário que a NMS envie algum tipo de solicitação para o agente para que o mesmo envie uma \emph{trap}, esse processo é assíncrono, ou seja, o agente informa ao gerente automaticamente quando algo acontece no dispositivo gerenciado.

\subsection{Agente}
Aplicação que é armazenada nos dispositivos gerenciados, e é responsável em coletar eventos e informar ao gerente o que acontece no dispositivo.
\par
Para coletar essas informações, o agente faz a varredura do dispositivo monitorado.
\par
Por exemplo, se colocar um agente para monitorar um servidor de banco de dados, e se o uso do processador desse servidor aumentar além do que foi configurado como quantidade limite de operação, o agente imediatamente envia uma \emph{trap} para gerente (NMS), onde será tratada, informando do evento que está acontecendo e a NMS pode, se estiver configurado, enviar um email avisando o responsável pelo servidor sobre o que está acontecendo com aquele servidor.

\subsection{MIB}
A MIB (\emph{Management Information Base}) é a base de informação do objeto gerenciado. É através da MIB que estão as informações coletadas pelos agentes.
\par
Quando uma NMS consulta informações de um dispositivo gerenciado, é na MIB que são consultadas essas informações.
\par
São definidas três tipos de MIBs padão, MIB II, MIB Experimental e MIB Privada.
\begin{enumerate}
\item[*] {\textbf{MIB II:}} Evolução da MIB I, descrita pela RFC 1213;
\item[*] {\textbf{MIB Experimental:}} Utilizada para desenvolvimento e testes de novas funcionalidades de gerenciamento;
\item[*] {\textbf{MIB Privada:}} Fornece informações especificas dos dispositivos gerenciados.
\end{enumerate}

\par
Todos os agentes contém a MIB II, pois essa MIB implementa os recursos necessários para monitoramento do protocolo TCP/IP, porém o agente pode implementar dos vários tipos de MIBs.
\par
Além dessas MIBs padão existem outras específicas para determinados tipos de serviços, por exemplo DNS Server MIB (RFC 1611) específica para servidores que executam serviços DNS.
\par
\textbf{Construção de MIB}
\par 
A construção da estrutura das MIBs é baseada em regras descritas através da SMI (\emph{Structure of Management Information}) - Estrutura de Informações de Gerenciamento. SMI define os objetos gerenciados.

\subsection{Versões do Protocolo}

%\begin{enumerate}
%\item RFC 3584\cite{rfc-3584}
%\item RFC 2119 - Best Current Practice
%\item RFC 2576 - Coexistencia entre SNMPv1, SNMPv2 e SNMPv3
%\item RFC 2578 - Define Estrutura SMIv2 - Structure Management Information
%\item RFC 2579 - Define common MIB
%\item RFC 2580 - Define declarações e requisitos definindo capacidade de agente e gerenciamento
%\item RFC 3416 - Define operações do protocolo
%\item RFC 3417 - Define mapeamento de transporte usado em cabo
%\item RFC 3418 - Gerenciamento de Informação basica entidades SNMP
%\item RFC 1901 - Definição experimental usando SNMPv2
%\item RFC 3410 - Define SNMPv3
%\item RFC 3411 - Define a Arquitetura Framework Gerenciamento SNMP
%\item RFC 3412 - Define o processo de mensagem
%\item RFC 3413 - Define varios aplicativos SNMP
%\item RFC 3414 - Define o modelo de Segurança baseado em usuario
%\item RFC 3415 - Define modelo de controle de acesso baseado em visualização
%\item RFC 1303 - Capacidade descrever agentes para um ou mais modulos MIB
%\item RFC 1908 - Coexistencia entre SNMPv1 e SNMPv2
%\item RFC 2089 - Agente SNMP
%
%\item RFC 1157 - Define SNMPv1
%\item RFC 1155 - Define Estrutura SMIv1 - Structure Management Information
%\item RFC 1212 - Define Mecanismo SMIv1 
%\item RFC 1215 - Convenção para definir Traps SMIv1
%\end{enumerate}

\section{Netflow / IPFIX} \label{sec:netflow}

\section{IEEE 802.1q}

\section{Intrusion Detection System - IDS}
\subsection{Snort}

\section{Nessus}


\chapter{A FERRAMENTA}

\section{Características}

\section{Componentes}

\section{Requisitos de Sistema}

\section{Requisitos de Hardware}


\chapter{IMPLEMENTAÇÃO E CONFIGURAÇÃO}

\section{PacketFence}

\section{Configuração Switch}


\chapter{CONCLUSÃO}

%\chapter{PROBLEMA A SER ABORDADO}
%O NAC pode ser utilizado em qualquer organização que queira implantar controles de quem pode ou não acessar a rede cabeada.
%\par
%Estes controles são baseados em políticas de acessos ao ambiente corporativo.
%
%\chapter{OBJETIVOS A SEREM ATINGIDOS}
%\section{GERAL}
%Demonstrar através de uma implantação em ambiente real uma maneira de gerenciar os acessos à rede cabeada.
%
%\section{OBJETIVOS ESPECÍFICOS}
%Pretende-se com este trabalho preencher a falha de segurança que os pontos de redes sem utilização apresentam, com o uso do DHCP na corporação, o atacante pode facilmente obter  as configurações referente a infra-estrutura da corporação, tais como: endereço ip, máscara de rede, gateway e dns, deixando facilmente acessível a infraestrutura de rede com a obtenção destas informações. Possibilitando ataques originados de dispositivos não autorizados na rede.
%\par
% É também parte do escopo deste trabalho apresentar uma tradução, dentro dos limites permissivos, para a língua portuguesa da documentação dos manuais do software PacketFence.
%
%\chapter{JUSTIFICATIVA DO PROJETO}
%Com a expansão das redes de computadores nos ambientes corporativos, se torna extremamente difícil o controle de quem pode ou não usar o cabeamento da organização. Portanto, uma pessoa com más intenções pode, se tiver conhecimento, utilizar dessa falha de segurança para roubar ou modificar dados sigilosos para os negócios da empresa, ou até mesmo ocasionar uma queda das operações eletrônicas da organização.
%\par
%Uma das intenções do trabalho é demonstrar como corrigir ou amenizar (fechar) essa falha com soluções open source existentes no mercado, pois o argumento dos gestores é que as soluções comerciais são caras, além de outra dificuldade bastante visível que seria a falta de mão de obra qualificada.
%
%\chapter{REFERÊNCIAS TEÓRICAS QUE O EMBASA}
%Serão utilizados os seguintes referenciais teóricos para o escopo deste trabalho e onde os seus respectivos conteúdo podem ser encontrados:
%
%\begin{enumerate}
%\item[-]Marcotte, Ludovic; Gehl, Dominik (2007-04-01). "PacketFence". Linux Journal. \url{http://www.linuxjournal.com/article/9551};
%\item[-]Inverse Inc., PacketFence Administration Guide. \url{http://www.packetfence.org/downloads/PacketFence/doc/PacketFence_Administration_Guide-2.2.1.pdf}. 03 de Ago 2011.
%\end{enumerate}
%
%\chapter{METODOLOGIA A SER UTILIZADA}
%Será realizada uma pesquisa exploratória em material escrito, como, por exemplo, artigo publicado pelos próprios desenvolvedores do sistema, para coleta e seleção de conteúdo pertinente ao escopo deste trabalho.
%\par
%Para a realização e desenvolvimento da parte prática, pode-se apontar a seguintes ações:
%\begin{enumerate}
%\item[-]{Configurar o servidor que fará as autenticações dos computadores e outros ativos de redes para utilizar a infraestrutura da empresa, sendo que esse servidor será o ponto central da rede, todas as conexões passarão por esse equipamento. Serão habilitadas as funcionalidades do software PacketFence (IDS, Nessus e Radius);}
%\item[-]{Faremos com que o servidor trabalhe juntamente com o switch para realizar o controle dos novos dispositivos que serão conectados à rede da empresa;}
%\item[-]{Tais configurações deverão ser realizadas com base na documentação original da ferramenta PacketFence, que será um dos principais documentos de pesquisa.}
%\end{enumerate}
%
%\chapter{CRONOGRAMA PARA A EXECUÇÃO DO TRABALHO DE TCC}
%\begin{tabular}{|p{4cm}|c|c|c|c|c|c|c|c|c|c|}
%\hline
%\textbf{Atividades}&		\textbf{Mar}&	\textbf{Abr}&		\textbf{Mai}&		\textbf{Jun}&		\textbf{Jul}&		\textbf{Ago}&	\textbf{Set}&		\textbf{Out}&	\textbf{Nov}\\
%\hline
%Elaboração do projeto&	&	&	&	X&	X&	X&	&	&	\\
%\hline
%Entrega do Projeto&		&	&	&	&	&	&	&	&	X\\
%\hline
%Pesquisa Bibliográfica&	&	&	&	X&	X&	X&	&		&\\
%\hline
%Coleta de Dados&		&	&	&	&	&	X&	&		&\\
%\hline
%Apresentação e discussão dos dados&	&	&	&	&	&	&	X&	&\\
%\hline
%Conclusão&	&	&	&	&	&	&	&	X&	\\
%\hline
%Entrega do TCC&	&	&	&	&	&	&	&	&	X\\
%\hline
%Defesa da Banca&	&	&	&	&	&	&	&	&	X\\
%\hline
%\end{tabular}

% --------------------------- BIBLIOGRAFIA ------------------------ %
%\chapter{REFEÊNCIAS BIBLIOGRÁFICAS}
% plain, abnt-alf, abnt-num, unsrt, abbrv, alpha,
\bibliographystyle{abnt-alf}
\bibliography{tcc}
%\addcontentsline{toc}{chapter}{Referências Bibliográficas}
%\singlespacing

\end{document}
